The JITJSS problem, introduced by \citet{baptiste2008lagrangian}, considers a set of $n$ jobs $J = \{J_1, J_2, \ldots, J_n\}$ to be processed on a set of $m$ machines $M = \{M_1, M_2, \ldots, M_m\}$.
Each job $J_i \in J$ consists of an ordered sequence of $m$ operations $O_i = \{o_{i1}, o_{i2}, \ldots, o_{im}\}$, where $o_{ik}$ represents the $k$-th operation of job $J_i$.
Each operation $o_{ik}$ must execute on a specified machine $M(o_{ik}) \in M$ for an uninterrupted processing time $p_{ik} > 0$.
We denote by $\mathcal{O}(M_u) = \{o_{ik} \mid M(o_{ik}) = M_u\}$ the set of all operations requiring machine $M_u \in M$.

Each operation $o_{ik}$ has a target due date $d_{ik} \ge 0$, an earliness penalty weight $\alpha_{ik} \ge 0$, and a tardiness penalty weight $\beta_{ik} \ge 0$.
Let $C_{ik}$ denote the completion time of operation $o_{ik}$.
The earliness $E_{ik}$ and tardiness $T_{ik}$ of operation $o_{ik}$ are defined as:
\begin{align}
E_{ik} &= \max(d_{ik} - C_{ik}, 0), \quad \forall J_i \in J, \, o_{ik} \in O_i, \label{eq:earliness_def} \\
T_{ik} &= \max(C_{ik} - d_{ik}, 0), \quad \forall J_i \in J, \, o_{ik} \in O_i. \label{eq:tardiness_def}
\end{align}

The objective of JITJSS is to determine a feasible schedule of operation completion times that minimizes the total weighted earliness and tardiness penalty $z$:
\begin{equation}
\min z = \sum_{i=1}^{n} \sum_{k=1}^{m} \left( \alpha_{ik} E_{ik} + \beta_{ik} T_{ik} \right) \label{eq:obj}
\end{equation}
subject to:
\begin{align}
&C_{i1} - p_{i1}     \ge 0, && \forall J_i \in J, \label{eq:non_negative_start} \\
&C_{ik} \le C_{i,k+1} - p_{i,k+1}, && \forall J_i \in J, \, k = 1, \ldots, m-1, \label{eq:job_precedence} \\
&C_{ik} \ge C_{jh} + p_{ik} \lor C_{jh} \ge C_{ik} + p_{jh}, && \forall M_u \in M, \, o_{ik} \neq o_{jh} \in \mathcal{O}(M_u), \label{eq:machine_disjunction} \\
&E_{ik} \ge d_{ik} - C_{ik}, && \forall J_i \in J, \, o_{ik} \in O_i, \label{eq:earliness_linear} \\
&T_{ik} \ge C_{ik} - d_{ik}, && \forall J_i \in J, \, o_{ik} \in O_i, \label{eq:tardiness_linear} \\
&E_{ik} \ge 0, \quad T_{ik} \ge 0, && \forall J_i \in J, \, o_{ik} \in O_i. \label{eq:non_negativity}
\end{align}

The objective function~\eqref{eq:obj} minimizes the total weighted penalty for both early and late operation completions.
Constraint~\eqref{eq:non_negative_start} ensures that no operation starts before time zero.
Constraint~\eqref{eq:job_precedence} enforces job precedence constraints, requiring that each operation $o_{ik}$ completes before its successor $o_{i,k+1}$ begins processing.
Constraint~\eqref{eq:machine_disjunction} represents disjunctive machine constraints, preventing any machine $M_u$ from processing more than one operation simultaneously.
Constraints~\eqref{eq:earliness_linear}--\eqref{eq:non_negativity} define the non-negative earliness and tardiness values for each operation.

\subsection{Sequence Representation}
\label{subsec:disjunctive_graph}

Following standard job shop scheduling literature \citep{pinedo2016}, we represent machine processing orders using a disjunctive graph $G = (V, X, Y)$.
The node set $V = \{s, t\} \cup \{o_{ik} \mid J_i \in J, k = 1, \ldots, m\}$ contains a dummy source node $s$, a dummy sink node $t$, and nodes representing all operations, with node weights equal to processing times $p_{ik}$.
The directed arc set $X$ represents fixed job precedence constraints, containing directed arcs $(s, o_{i1})$, $(o_{im}, t)$, and $(o_{ik}, o_{i,k+1})$ for $k = 1, \ldots, m-1$, which enforce the fixed sequence of operations within each job $J_i$.
The undirected edge set $Y$ represents potential machine conflicts, containing an undirected edge $\{o_{ik}, o_{jh}\}$ for every pair of operations requiring the same machine $M_u$.

A solution sequence $\pi$ specifies a processing order for each machine, which corresponds mathematically to directing every undirected edge $\{o_{ik}, o_{jh}\} \in Y$ as $(o_{ik} \to o_{jh})$ if $o_{ik}$ precedes $o_{jh}$ on machine $M_u$.
This complete selection yields a fully directed graph $G(\pi) = (V, X \cup Y(\pi))$, where any directed path between two operations indicates an enforced temporal precedence constraint.
A sequence $\pi$ represents a physically feasible schedule if and only if the directed graph $G(\pi)$ contains no directed cycles, as a cycle would imply a deadlocked precedence loop where an operation must precede itself \citep{pinedo2016}.

In traditional makespan minimization, setting operation start times as early as possible on a given acyclic graph $G(\pi)$ yields an optimal schedule, which researchers compute efficiently in $O(|V| + |X|)$ time via critical path calculation.
In contrast, JITJSS penalizes early completions, meaning that scheduling operations as early as possible is generally non-optimal.
For a fixed feasible operation sequence $\pi$, determining the optimal completion times $C_{ik}$ requires solving a timing sub-problem to balance earliness and tardiness penalties.
Metaheuristic methods for JITJSS typically decompose the solution process into two stages: exploring the space of operation sequences $\pi$ via local search or metaheuristics, and solving the timing sub-problem for each sequence using mathematical programming or linear programming sub-solvers \citep{dos2010combination, wang2014variable, ahmadian2021meta}.
